\section{Propuesta para datos semiestructurados, no estructurados y vectoriales}

Cada fragmento documental es una unidad vectorizada. \texttt{embedding}
referencia al fragmento y al modelo. La versión, el documento, la vigencia y la
clase se obtienen mediante relaciones. La vista
\texttt{fragmento\_recuperable} reúne esos datos sin copiarlos en el vector, y
RLS aplica los permisos del perfil efectivo sobre las tablas base.

Los vectores sintéticos tienen 32 dimensiones, están normalizados y forman
vecindarios temáticos deterministas. La búsqueda usa distancia coseno
(\texttt{<=>}), ordena de menor a mayor y desempata por
\texttt{fragmento\_id}. El modelo activo cubre todos los fragmentos. Cuatro
también conservan embeddings históricos, excluidos de la recuperación vigente.

El índice HNSW usa \texttt{vector\_cosine\_ops}. La muestra es demasiado
pequeña para evaluar rendimiento: PostgreSQL elige una exploración secuencial
(\texttt{Seq Scan}) porque recorrer 32 vectores activos cuesta menos que usar el
índice. Una ejecución diagnóstica con \texttt{enable\_seqscan = off} confirma
que HNSW puede resolver el patrón top-k, pero ese plan no se fuerza.

La recuperación debe evitar cuatro errores: usar versiones obsoletas, exponer
clases no autorizadas, mezclar modelos y responder sin fuente. La vista filtra
actividad, publicación, vigencia y modelo. RLS aplica el permiso antes del
\texttt{ORDER BY} y del \texttt{LIMIT}; la unicidad evita duplicados; y la
evidencia guarda el \texttt{embedding\_id} exacto. Así una respuesta histórica
sigue siendo explicable aunque cambie la versión vigente.
